\section{Introduction}
\label{sec:rewrite-introduction}

Advances in artificial intelligence raise both hopes and fears. One hope is that increasingly capable AI could generate unprecedented economic prosperity, with workers benefiting from higher wages. One fear is that the AI industry could capture a growing share of that prosperity while workers who depend on labor income receive a declining share. In this paper, I show how both outcomes can coexist: when AI and labor are substitutes, AI can permanently raise wage growth. Yet in the long run, output per worker grows even faster than wages, while AI-industry revenue grows faster than total labor income. Labor’s income share therefore converges to zero.

To show this, I extend the Ramsey--Cass--Koopmans model to include AI services as an additional production input. Households choose saving, and competitive firms produce the final good using capital, effective labor, and AI services. A monopolistic AI developer supplies these services and invests in recursive AI self-improvement (RSI) to improve AI efficiency---the quantity of AI services produced per unit of compute. As in \citet{romer1990}, expected monopoly profits provide incentives for research investment. %Together, the elasticity of substitution between AI production services and labor and the maximum attainable level of AI efficiency determine whether human labor remains a long-run production bottleneck or whether AI accelerates economic growth.
By introducing a monopolistic AI developer, I can distinguish how the gains from AI are distributed through wages, returns to physical capital, and AI-industry profits. 

%This fear, however, does not require real wages to decline: workers may receive a smaller share of a much larger economy even as their wages rise. Can unprecedented prosperity coexist with a vanishing labor income share? Under what conditions could the world economy follow such a path? How long could it take before the transition became visible, and what economic signals would reveal it? What would happen to wages and labor's income share along the way?

%This question is becoming more pressing as AI increasingly contributes to its own development. Recursive self-improvement (RSI) is a feedback process in which better AI helps develop still better AI. OpenAI \citep{openairesearchacceleration2026} and Anthropic \citep{favaroclark2026} report that AI systems already accelerate parts of their research. OpenAI's chief scientist already stated that the pace of progress ``could be sustained into recursive self-improvement'' \citep{pachocki2026}. Anthropic nevertheless cautions that, regarding AI systems that autonomously design and develop their successors, ``We are not there yet'' \citep{favaroclark2026}. These reports motivate studying the economic consequences of autonomous AI research without assuming when it will become possible.


The model combines two feedback loops related to those in \citet{davidsonetal2026} and \citet{jonestonetti2026}. A technological loop operates through RSI: greater AI efficiency produces more AI research services from a given amount of compute devoted to research and thereby helps create still better AI. An economic loop runs from AI efficiency to output and back to research: greater AI efficiency raises production and expands the resources available to fund further research.

With substitution and sufficiently strong research returns, the developer's discounted net profit can become unbounded if AI efficiency is unrestricted. To rule this out, I first study an economy with an upper bound on AI efficiency and examine how equilibria change as the bound rises. Together, the elasticity of substitution between AI production services and labor and the maximum attainable level of AI efficiency determine whether human labor remains a long-run production bottleneck or whether AI sustains permanently faster economic growth. %I focus on two long-run equilibrium regimes, distinguished by whether AI produces a temporary or permanent increase in growth. In the labor-bottleneck regime, AI improvements can raise growth temporarily, but human labor ultimately constrains growth. In the AI-dominated regime, capital and AI services sustain permanently faster growth.

I first characterize the equilibria in which labor remains a long-run production bottleneck. %Under complementarity, I establish their existence when the upper bound on AI efficiency exceeds an elasticity-dependent threshold. At unit elasticity, such equilibria exist for any upper bound. With substitution, they exist when the upper bound lies below the corresponding threshold. 
Along these equilibria, output per worker and wages eventually grow at the exogenous rate of labor-augmenting technological progress, and labor's income share converges to a positive value. This is the result when there is complementary between labor and AI production services or when the upper bound on AI efficiency is sufficiently low.

An AI-dominated long-run regime then becomes possible when the elasticity of substitution between AI production services and labor exceeds one and the upper bound is sufficiently high. As in \citet{jonesmanuelli1990}, capital accumulation can sustain growth when its limiting return remains sufficiently high. In this regime, capital and AI production services eventually grow faster than effective labor. This does not require AI efficiency to keep improving: capital accumulation and expanding compute use sustain growth even as AI efficiency approaches its upper bound. The long-run growth rates of output per worker and wages, as well as the interest rate, exceed their labor-bottleneck values. Because output per worker grows faster than wages, labor's income share converges to zero. %Within the AI-dominated regime, long-run growth and the interest rate increase with the maximum attainable level of AI efficiency.
Substitution thus permits a permanent acceleration of wages even as labor's share vanishes: capital and AI services expand beyond the pace of effective labor, raising each worker's marginal product. Total labor income therefore grows faster than in the labor-bottleneck regime, but in the long run it grows more slowly than capital income and AI-industry revenue.

%\Needspace{3\baselineskip}
%The distinction between regimes concerns the persistence of growth, not whether AI produces economic gains. An economy can benefit from more efficient AI and sustain a profitable AI industry without permanently increasing output-per-worker growth. In the AI-dominated regime, a permanent increase in output growth translates into wage growth: the wage-growth premium above the exogenous rate of labour-augmenting technological progress equals the output-per-worker growth premium divided by the elasticity of substitution. Workers therefore receive higher real wages, but these wages do not keep pace with the economy as a whole.

I also characterize a competitive economy in which AI models and any improvements are freely available to all suppliers. Without a private return to research, AI efficiency remains at its initial level. When AI and labor are substitutes, a sufficiently high initial efficiency nevertheless sustains AI-dominated growth through capital accumulation and expanding compute use. Neither monopoly nor continuing RSI is therefore necessary to sustain this growth regime.

The growth comparison between monopoly and competition depends on whether monopoly-financed efficiency gains compensate for the reduction in AI use caused by monopoly pricing. The simulations show that research can allow the monopoly economy to reach AI-dominated growth from an initial efficiency too low to sustain it under competition. Conversely, competition can sustain AI-dominated growth from initial states for which monopoly-led research improves efficiency but leaves the economy labor-bottlenecked. Even when both economies approach AI-dominated growth, monopoly delivers faster long-run growth only if the research gains are large enough to offset its markup. More efficient AI therefore need not imply faster long-run growth.

I also examine the effect of removing the upper bound on AI efficiency. With complementarity, long-run average output-per-worker growth cannot exceed the exogenous rate of labor-augmenting technological progress. At unit elasticity, I derive conditions for a balanced-growth equilibrium in which output per worker grows faster than this exogenous rate. Wages grow at the same rate as output per worker, preserving a constant labor income share. With substitution and sufficiently strong research returns, the developer can make discounted net profit arbitrarily large even over a fixed finite horizon, ruling out the existence of an equilibrium.

I use numerical simulations to illustrate the effects of introducing RSI under monopoly and of replacing competition with monopoly. First, I simulate an unanticipated activation of RSI in economies that already use AI services. Depending on research productivity, output-per-worker growth can display a moderate increase or a noticeable jump in both long-run regimes. Eventually, output-per-worker growth converges to the rate of labor-augmenting technological progress in the labor-bottleneck regime and to a permanently higher rate in the AI-dominated regime. With lower research productivity, initial output-growth trajectories can remain similar across economies approaching different long-run regimes. Second, I simulate granting exclusive AI rights in economies that initially supply AI services competitively, separating the immediate reduction in AI use from subsequent research-driven efficiency gains. Monopoly-financed research can move an economy from labor-bottleneck to AI-dominated growth. Conversely, monopoly pricing can leave an economy labor-bottlenecked even when competition would sustain AI-dominated growth with less efficient AI. %Higher research productivity induces greater initial research expenditure and a larger reduction in capital accumulation, even as consumption rises. Current developer profits fall more sharply because more revenue is devoted to research. Lower research productivity produces a milder initial adjustment, and economies approaching different long-run regimes can initially display similar growth. %Early aggregate growth therefore provide incomplete information about the economy's eventual trajectory.

This paper has eight sections, including the present introduction. Section~\ref{sec:rewrite-literature} relates the paper to the literature. Section~\ref{sec:rewrite-model} presents the model and defines equilibrium. Section~\ref{sec:rewrite-regimes} establishes sufficient conditions for equilibrium existence with an upper bound on AI efficiency and characterizes the associated long-run growth and distributional outcomes. Section~\ref{sec:rewrite-competition} characterizes competitive AI provision and compares its growth outcomes and research incentives with those under monopoly. Section~\ref{sec:rewrite-uncapped} examines the three substitution regimes without an upper bound on AI efficiency. Section~\ref{sec:rewrite-quantitative} presents numerical simulations. Section~\ref{sec:rewrite-conclusion} concludes. The appendix contains the proofs and describes the algorithm used for the simulations.
